\section{Three linear systems for the Bessel comparison}

% BEGIN L90I1
\paragraph{L90I1. Computing the spaces of intertwiners.}
On a datum of B90T1 put $A_u=H(T_u)$ and $C_u=H(S_u)$, acting
on the common row space $E$ of dimension $d$. Thus $A_u$ and $C_u$
are self-adjoint for the row inner product. The bilinearity needed to
form these operators holds on every native datum. The two linearity
axioms for $B_u$ are fields of $D7$. The source proves the corresponding
two finite-sum identities for $\mathrm{besS}_u$ with unconjugated
coefficients in both arguments.

For $P,Q\in\{A,C\}$ consider the complex linear map
\[
  D_{PQ}:\operatorname{End}_{\mathbb C}(E)\longrightarrow
           \bigoplus_{u}\operatorname{End}_{\mathbb C}(E),
  \qquad X\longmapsto(Q_uX-XP_u)_u.
\]
Write $h_{PQ}=\dim_{\mathbb C}\ker D_{PQ}$. The column spaces and
their metrics have already entered through the Grams in B90T1;
$X$ acts only on the row space. For example, using a rational row
basis with Gram $H$ gives the matrices of $A_u,C_u$ in that basis.
They satisfy $P_u^*H=HP_u$. The linear equations in the displayed
formula are valid in these same coordinates, and a simultaneous
change of row basis preserves their nullities.

There are finitely many labels by the $\mathrm{wordFin}$ field of
$D7$. Stack the blocks $I_d\otimes Q_u-P_u^t\otimes I_d$ over all
labels, using column-by-column vectorization. Then $h_{PQ}$ is
$d^2$ minus the rank of this matrix, which has $d^2$ columns and
$|\mathrm{Word}|d^2$ rows. The transpose in the block formula is
ordinary transpose, since the equations are complex linear in $X$.
% END L90I1

% BEGIN L90I2
\paragraph{L90I2. A sum of squares detects simultaneous unitary equivalence.}
The three nullities $h_{AA},h_{CC},h_{AC}$ give a nonnegative
integer
\[
                     \Delta=h_{AA}+h_{CC}-2h_{AC}.
\]
We claim that $\Delta=0$ precisely when there is a unitary $V$ of
$E$ with $C_uV=VA_u$ for every $u$.

To prove this, decompose each row space orthogonally into subspaces
irreducible under its family of operators. Such a decomposition
exists by induction on dimension, since self-adjointness makes the
orthogonal complement of each common invariant subspace invariant.
Group these summands from both families by unitary equivalence
with all labels preserved. For class $j$, write $m_j$ and $n_j$
for its multiplicities on the $A$ and $C$ sides respectively, taking
zero for an absent class.

An intertwiner between irreducibles is zero or invertible by its
invariant kernel and image. Its adjoint intertwines in the reverse
direction. An endomorphism of an irreducible is scalar, since the
eigenspace for any of its complex eigenvalues is nonzero and
invariant. Hence for a nonzero intertwiner $X$ the operator $X^*X$
is a positive scalar times the identity. Thus the intertwiner
space between two irreducibles is zero for distinct classes and a
complex line for equal classes. Taking blocks of arbitrary
intertwiners between the orthogonal sums yields
\[
 h_{AA}=\sum_j m_j^2,\qquad h_{CC}=\sum_j n_j^2,\qquad
 h_{AC}=\sum_j m_jn_j,\qquad
 \Delta=\sum_j(m_j-n_j)^2 .
\]
If the multiplicities agree, match the summands class by class
using unitary intertwiners and take their orthogonal sum. Conversely
a common unitary gives linear isomorphisms among the three
intertwiner spaces by composition, so $\Delta=0$.
% END L90I2

% BEGIN L90I3
\paragraph{L90I3. The remaining rank identity on the native family.}
By B90T2 the unitary in L90I2 gives the original wordwise column
unitaries and nonzero scalars. Hence on any datum for which the
native rationality clause holds, including the selected range of
R88L3, the entire Bessel identity is equivalent to
\[
      2\operatorname{rank}D_{AC}
           =\operatorname{rank}D_{AA}+\operatorname{rank}D_{CC}.
\]
For an arbitrary datum one must also establish the rationality
clause. A simultaneous intertwining unitary preserves every zero
label, so the rank criterion already detects differing zero
supports. If both families vanish identically, the three kernels
all have dimension $d^2$ and the unitary comparison holds. This includes
$d=0$, with the unique unitary on the zero space. For $d=1$ every
normalized Gram is zero or one, and agreement of the zero supports
suffices. For general $d$, the displayed equation gives an exact
test using only the three ranks. Its validity for the selected
arithmetic and spectral operators still requires a mathematical
identity for that family.
% END L90I3
